\documentclass[11pt,a4paper]{article}

% ---------- packages ----------
\usepackage[margin=1in]{geometry}
\usepackage[utf8]{inputenc}
\usepackage[T1]{fontenc}
\usepackage{amsmath,amssymb,amsthm}
\usepackage{graphicx}
\usepackage{xcolor}
\usepackage{hyperref}
\usepackage{booktabs}
\usepackage{array}
\usepackage{multirow}
\usepackage{float}
\usepackage{caption}
\usepackage{subcaption}
\usepackage{enumitem}
\usepackage{fancyhdr}
\usepackage{lastpage}
\usepackage{natbib}

% ---------- theorem environments ----------
\theoremstyle{plain}
\newtheorem{theorem}{Theorem}[section]
\newtheorem{lemma}[theorem]{Lemma}
\newtheorem{corollary}[theorem]{Corollary}
\newtheorem{conjecture}[theorem]{Conjecture}

\theoremstyle{definition}
\newtheorem{definition}[theorem]{Definition}
\newtheorem{remark}[theorem]{Remark}
\newtheorem{example}[theorem]{Example}

% ---------- page style ----------
\pagestyle{fancy}
\fancyhf{}
\lhead{Du}
\chead{Missing-Center Solutions in the No-Three-In-Line Problem}
\rhead{\thepage}
\renewcommand{\headrulewidth}{0.4pt}

% ---------- footer ----------
\fancypagestyle{firstpage}{
  \fancyhf{}
  \lfoot{\small \textcolor{gray}{This is a preprint. Submitted for publication.}}
  \rfoot{\small \textcolor{gray}{\thepage}}
}

% ---------- title ----------
\title{\vspace{-1.5cm}\textbf{Missing-Center Solutions in the\\No-Three-In-Line Problem:\\
A Computational and Structural Analysis}}

\author{Junrong Du\\
\small Independent Researcher\\
\small \texttt{dududu9738@gmail.com}\\[4pt]
\small MSC2020: 05B99, 68R05\\
\small Keywords: No-Three-In-Line, combinatorial geometry, exhaustive search,\\
\small \hspace*{1.85cm}circumcenter invariant, C$_4$ symmetry, distance ring}

\date{\small \today}

\begin{document}

\maketitle
\thispagestyle{firstpage}

% ============================================================
\begin{abstract}
The No-Three-In-Line problem asks for the maximum number of points that can be placed on an $n\times n$ grid with no three collinear. This paper studies the geometric structure of maximal $2n$-point configurations through a novel invariant: whether the grid center serves as a circumcenter of any triple in the configuration---an invariant we show reduces to an exact integer distance condition.

Our main results are three theorems about symmetry-imposed structure in extremal configurations:

\begin{enumerate}[nosep]
\item \textbf{C$_4$ Theorem}: Every $2n$-point solution with $90^\circ$ rotational symmetry has the grid center as a circumcenter. The proof is a one-line distance invariance under rotation.
\item \textbf{Row-Degree Theorem}: For even $n=2m$, any C$_4$-symmetric solution corresponds bijectively to a $2$-regular graph on $m$ vertices, reducing the existence problem from geometry to graph theory.
\item \textbf{C$_2$ Direction Theorem}: In $180^\circ$-symmetric solutions, the $n$ antipodal pairs must occupy $n$ distinct directions from the center; collinearity between pairs occurs iff they share a direction. This reveals an unexpected phase transition at $n=31$ where $44{,}828$ solutions exist at $n=29$ yet zero at $n=31$.
\end{enumerate}

These structural constraints explain why missing-center solutions---those lacking the center as a circumcenter---are restricted to specific symmetry classes. We provide exhaustive missing-center counts for $n=2\dots13$ and $D_4$-inequivalent analysis through $n=53$, and identify $n=71$ as the sole remaining open case $\le 72$. The evidence strongly supports $D(n)=2n$ for all even $n$.
\end{abstract}

% ============================================================
\section{Introduction}
\label{sec:intro}

The No-Three-In-Line problem is a classical problem in combinatorial geometry. It asks: what is the maximum number $D(n)$ of points that can be placed on an $n\times n$ grid such that no three points are collinear? The problem has been studied since at least the early 20th century~\cite{erdos1935}, with significant computational progress in recent decades~\cite{guy1994,flammenkamp}.

It is known that $D(n) \le 2n$ (by the pigeonhole principle on rows) and it is conjectured that $D(n) = 2n$ for all $n$ except possibly a finite set. For many years, the precise values of $D(n)$ were known only up to $n=52$ (the classical bound from exhaustive search before the 2026 SAT-solver advances of Heule). In 2026, Heule~\cite{heule2026} dramatically advanced the state of the art using SAT solvers, establishing $D(n)=2n$ for all $n\le 72$ with the sole exception of $n=71$.

This paper approaches the problem from a new perspective. Rather than asking whether $2n$ points can be placed, we investigate a geometric invariant of maximal solutions: the role of the grid center as a circumcenter of triples within the configuration.

\begin{definition}
A \emph{maximal solution} is a set of exactly $2n$ grid points with no three collinear. The grid center is $C = \bigl(\frac{n-1}{2},\frac{n-1}{2}\bigr)$. A solution is \emph{missing-center} (or \emph{center-free}) if no triple of its points has $C$ as its circumcenter.
\end{definition}

The detection of center-circumcenter triples reduces to an elegant integer criterion:

\begin{lemma}
\label{lem:distance}
The grid center $C$ is a circumcenter of three grid points iff the three points share the same squared Euclidean distance from $C$.
\end{lemma}

\begin{proof}
Points on a circle centered at $C$ are exactly those at a fixed distance from $C$. The squared distance is $d(i,j) = (2i-(n-1))^2 + (2j-(n-1))^2$, which is an integer. Conversely, if $C$ is the circumcenter of three points, they are equidistant from $C$, hence share the same $d$ value. The equivalence is exact --- no floating-point approximation is involved.
\end{proof}

\noindent This integer-criterion allows us to detect missing-center solutions using exact arithmetic.

\paragraph{Why the missing-center invariant?} Although the center being a circumcenter does not directly affect collinearity or point count, it provides a new classification axis for maximal solutions. Our three symmetry theorems---C$_4$, Row-Degree, and C$_2$---each reveal how the choice of symmetry class imposes structural constraints on the configuration, and the missing-center property is one of the clearest manifestations of these constraints. Configurations that lack the center as a circumcenter are precisely those that avoid placing four points on any distance ring, a highly restrictive condition that selects for specific symmetry classes (rot2 and identity). Understanding this restriction helps explain the observed phase transitions and may inform search strategies for the remaining open case $n=71$.

\paragraph{Main contributions.}
\begin{enumerate}[nosep]
\item \textbf{C$_4$ Theorem} (Theorem~\ref{thm:c4}): $90^\circ$ rotational symmetry forces the center to be a circumcenter.
\item \textbf{Row-Degree Theorem} (Theorem~\ref{thm:rowdeg}): C$_4$ solutions correspond to $2$-regular graphs, reducing geometry to graph theory.
\item \textbf{C$_2$ Direction Theorem} (Theorem~\ref{thm:c2}): $180^\circ$ symmetry forces $n$ distinct directions and a sharp phase transition at $n=31$.
\item Complete missing-center counts for $n=2\dots13$ (full search) and $D_4$-inequivalent analysis through $n=53$.
\item Characterization of the even-$n$ threshold at $n=12$ and identification of $n=71$ as the sole remaining open case $\le 72$.
\end{enumerate}

% ============================================================

\section{The C$_4$ Theorem}
\label{sec:c4}

A solution has \emph{C$_4$ symmetry} if it is invariant under $90^\circ$ rotation about the grid center. The rotation is $R(i,j) = (n-1-j,\, i)$.

\begin{theorem}[C$_4$ Theorem]
\label{thm:c4}
Let $S$ be a $2n$-point No-Three-In-Line solution with C$_4$ rotational symmetry. Then the grid center is a circumcenter of some triple in $S$.
\end{theorem}

\begin{proof}
Consider the squared distance $d(i,j) = (2i-(n-1))^2 + (2j-(n-1))^2$ from the center $C$. A direct computation shows $d$ is invariant under $R$: $d(R(i,j)) = d(i,j)$. The orbits of $R$ partition the grid points into sets of size 4 for even $n$, and size 4 plus one fixed point (the center) for odd $n$.

\textbf{Case 1: $n$ even.} Here $n=2m$. Then $2n=4m$, and every point belongs to a 4-orbit. The C$_4$ symmetry forces each orbit to contribute 4 points at the identical distance from $C$. Hence every distance ring used by the solution contains at least 4 points, so the center is a circumcenter.

\textbf{Case 2: $n$ odd.} Here $n=2m+1$. Then $2n=4m+2\equiv2\pmod{4}$. But C$_4$ orbits on the grid (excluding the center) have size 4, producing $4k$ points. Including the center gives $4k+1$ points. Neither $4k$ nor $4k+1$ equals $4m+2$. Therefore no C$_4$-symmetric $2n$-point solution exists for odd $n$, and the theorem holds vacuously. $\blacksquare$

\noindent\textbf{Remark.} The theorem is substantively non-trivial only for even $n$, where C$_4$-symmetric $2n$-point solutions do exist. For odd $n$, the theorem is vacuously true because no such solutions can exist by parity. This distinction matters: when we later discuss Heule's rct4 solutions for odd $n=65,67,69$, those belong to the rct4 symmetry class (which includes diagonal reflection), not the pure C$_4$ class. The rct4 class avoids the parity obstruction because its fundamental domain has size 2 rather than 4.
\end{proof}

\begin{corollary}
Missing-center solutions cannot have C$_4$ symmetry.
\end{corollary}

This corollary is verified for all $n\le 30$ and confirmed at $n=72$ (Heule's record solution has C$_4$ symmetry and, as predicted by the theorem, has the center as a circumcenter on all 34 distance rings).

\begin{theorem}[Row-Degree Equivalence for C$_4$]
\label{thm:rowdeg}
Let $n=2m$ be even and let $S$ be a C$_4$-symmetric $2n$-point No-Three-In-Line solution. Represent $S$ by the set of $m$ fundamental-domain cells $\{(i_1,j_1),\dots,(i_m,j_m)\}$ with $i_k,j_k\in[0,m-1]$ that generate the $m$ C$_4$ orbits of $S$. For each $k\in[0,m-1]$, define
\[
\deg(k) = \bigl|\{\text{selected orbits }(i,j): i=k\}\bigr| + \bigl|\{\text{selected orbits }(i,j): j=k\}\bigr|.
\]
Then $\deg(k)=2$ for all $k\in[0,m-1]$.
\end{theorem}

\begin{proof}
The C$_4$ rotation $R(i,j) = (n-1-j,i)$ acts on the $n\times n$ grid. An orbit $\{(i,j), R(i,j), R^2(i,j), R^3(i,j)\}$ consists of four points:
\[
(i,j),\quad (j,n-1-i),\quad (n-1-i,n-1-j),\quad (n-1-j,i).
\]
The rows occupied by this orbit are $\{i,\;j,\;n-1-i,\;n-1-j\}$; each of these four rows receives exactly one point from the orbit.

Now fix a row $r\in[0,n-1]$. A selected orbit $(i,j)$ contributes a point to row $r$ in exactly two cases:
\begin{itemize}
\item $i=r$ or $j=r$ (the point lands in row $r$ via the identity or the $270^\circ$ rotation); or
\item $n-1-i=r$ or $n-1-j=r$, i.e.\ $i=n-1-r$ or $j=n-1-r$ (the point lands in row $r$ via the $180^\circ$ or $90^\circ$ rotation).
\end{itemize}
Hence the number of points in row $r$ equals $\deg(r)$ when $r<m$, and equals $\deg(n-1-r)$ when $r\ge m$. Since every row must contain exactly $2$ points, we have $\deg(k)=2$ for all $k\in[0,m-1]$.
\end{proof}

\begin{corollary}[Graph Representation]
\label{cor:c4graph}
A C$_4$-symmetric $2n$-point solution on an even $n\times n$ grid corresponds bijectively to a $2$-regular graph (allowing loops and multiple edges) on $m=n/2$ vertices, where each selected fundamental-domain cell $(i,j)$ becomes an edge between vertices $i$ and $j$.
\end{corollary}

\begin{proof}
By Theorem~\ref{thm:rowdeg}, each vertex $k$ has $\deg(k)=2$, which is exactly the definition of a $2$-regular graph. An edge $(i,j)$ is a loop when $i=j$, contributing $2$ to $\deg(i)$ from a single edge, which is permitted in a $2$-regular multigraph.
\end{proof}

\noindent\textbf{Computational verification.} The theorem is satisfied by all $33{,}534$ known C$_4$ solutions in the Flammenkamp database ($n=12\dots56$) and by Heule's $n=72$ solution. Every distinct rot4 solution count per $n$ in Table~\ref{tab:c4counts} sums to $33{,}534$.

\noindent\textbf{Consequence for $D(n)$.} Theorem~\ref{thm:rowdeg} reduces the search for C$_4$-symmetric solutions from a geometric problem to a purely graph-theoretic one: find a $2$-regular graph on $m$ vertices whose $4m$ realized points contain no three collinear. Empirical evidence shows that a full $m$-cycle and the $(1,m-1)$ pattern are valid for all $m\in[10,36]$ (i.e., $n=20$ to $n=72$) for which detailed data exists, with isolated exceptions at $m=6,9$ correlated with divisibility by $3$. A complete proof that a collinearity-free $2$-regular graph exists for every $m$ would establish $D(2m)=2m$ for all $m$; the equivalence is a reformulation, not yet a constructive existence theorem.

% ============================================================
\section{The C$_2$ Direction Theorem}
\label{sec:c2}

A solution has \emph{C$_2$ symmetry} (also called $180^\circ$ rotational symmetry) if it is invariant under the rotation $R_{180}(i,j) = (n-1-i,\, n-1-j)$. This is the second-most restrictive symmetry class after C$_4$, and it governs a large fraction of known missing-center solutions.

\begin{theorem}[C$_2$ Direction Theorem]
\label{thm:c2}
Let $S$ be a $2n$-point No-Three-In-Line solution with $180^\circ$ rotational symmetry on an odd $n=2m+1$ grid. The $2n$ points partition into $n$ antipodal pairs $\{P, R_{180}(P)\}$. For two such pairs $p_1 = \{(i_1,j_1),(2m-i_1,2m-j_1)\}$ and $p_2 = \{(i_2,j_2),(2m-i_2,2m-j_2)\}$, the four points are collinear iff
\[
(i_1-m)(j_2-m) = (i_2-m)(j_1-m).
\]
\end{theorem}

\begin{proof}
Center coordinates at $C=(m,m)$. For a point $P=(i,j)$, its antipode is $R_{180}(P) = (2m-i,2m-j)$. The line through $P$ and $R_{180}(P)$ passes through $C$, with reduced direction vector $(i-m, j-m)$.

Now consider two pairs $p_1, p_2$. The four points are $\{P_1, R_{180}(P_1), P_2, R_{180}(P_2)\}$. If any triple among these four is collinear, by $180^\circ$ symmetry the four points are collinear along a line through $C$. For this to happen, the two direction vectors $(i_1-m, j_1-m)$ and $(i_2-m, j_2-m)$ must be collinear, i.e., their cross product vanishes:
\[
(i_1-m)(j_2-m) - (i_2-m)(j_1-m) = 0,
\]
which is exactly the stated condition.
\end{proof}

\begin{corollary}
\label{cor:c2dirs}
In any C$_2$-symmetric $2n$-point solution, the $n$ antipodal pairs must occupy $n$ distinct reduced directions from the grid center.
\end{corollary}

\begin{proof}
If two pairs share the same direction, the four points are collinear by Theorem~\ref{thm:c2}, violating the No-Three-In-Line condition.
\end{proof}

\paragraph{The rot2 phase transition at $n=31$.}
The C$_2$ Direction Theorem transforms the existence problem into a direction-packing problem: choose $n$ distinct reduced directions from the $\sim(n^2-1)/2$ available such that each row receives exactly $2$ points. The direction pool is always abundant---at $n=31$, there are $480$ candidate directions for $31$ needed, a $15\times$ surplus---so the bottleneck is not capacity but the interaction between the direction-uniqueness constraint and the collinearity constraints among off-center lines.

The data reveals a sharp satisfiability collapse:

\begin{table}[H]
\centering
\caption{SAT phase transition for C$_2$ (rot2) symmetry.}
\label{tab:rot2transition}
\begin{tabular}{crrc}
\toprule
$n$ & Solutions & $\binom{2n}{3}$ & $\binom{2n}{3} / ((n^2-1)/2)$ \\
\midrule
23 & 3,967 & 15,180 & 57.5 \\
25 & 8,980 & 19,600 & 62.8 \\
27 & 17,332 & 24,804 & 68.1 \\
29 & 44,828 & 30,856 & 73.5 \\
31 & \textbf{0} & 37,820 & \textbf{78.8} \\
33 & 1$^\dag$ & 45,760 & 84.1 \\
\bottomrule
\multicolumn{5}{l}{\footnotesize $^\dag$ One rot2 solution known at $n=33$ (Flammenkamp database).}
\end{tabular}
\end{table}

At $n=29$, there are $44{,}828$ C$_2$-symmetric solutions; at $n=31$, zero. A structured combinatorial threshold exists at a density of approximately $74$ potential collinear triples per available C$_2$ pair. The constraint graph is highly structured rather than random, so the analogy to random $k$-SAT is only qualitative. Finding the precise obstruction---whether it is a parity condition, a Hall-type matching failure, or a more subtle invariant---remains an open problem.

\paragraph{Relation to rct4.}
The rct4 symmetry class (which augments C$_4$ with diagonal reflection) has fundamental domain size $2$ rather than $4$, giving it orbits of size $4$ on odd grids. This class therefore inherits both the C$_4$ obstruction (center as circumcenter) and the C$_2$-type direction structure. The rct4 class is the only known symmetry class that supports $2n$-point solutions for odd $n\ge 33$, as demonstrated by Heule's $n=65,67,69$ results~\cite{heule2026}.

% ============================================================
\label{sec:results}

\subsection{Full Enumeration ($n=2$ to $n=13$)}

Using a local exhaustive backtracking search with the 2-per-row constraint, we obtained complete counts:

\begin{table}[H]
\centering
\caption{Full enumeration results: total and missing-center solutions.}
\label{tab:full}
\begin{tabular}{crrrr}
\toprule
$n$ & Type & Total & With Center & Missing \\
\midrule
2 & Even & 1 & 1 & 0\\
3 & Odd & 2 & 2 & 0\\
4 & Even & 11 & 11 & 0\\
5 & Odd & 32 & 28 & 4\\
6 & Even & 50 & 50 & 0\\
7 & Odd & 132 & 128 & 4\\
8 & Even & 380 & 380 & 0\\
9 & Odd & 368 & 360 & 8\\
10 & Even & 1135 & 1135 & 0\\
11 & Odd & 1120 & 1084 & 36\\
12 & Even & 4348 & 4296 & 52\\
13 & Odd & 3622 & 3330 & 292\\
\bottomrule
\end{tabular}
\end{table}

\subsection{$D_4$-Inequivalent Analysis ($n=7$ to $n=30$)}

We extended the analysis using the Flammenkamp database~\cite{flammenkamp} of $D_4$-inequivalent solutions, organized by symmetry class. This provides data for larger $n$ where full enumeration is infeasible.

\begin{table}[H]
\centering
\caption{$D_4$-inequivalent missing-center counts from the Flammenkamp database (full symmetry classes through $n=45$; rct4-only entries for $n=47,49,51,53$ from mvr/no-three-in-line). Even $n=46,48,50,52$ are omitted because only rot4 counts are catalogued at those sizes (see Table~\ref{tab:c4counts}).}
\label{tab:d4}
\footnotesize
\begin{tabular}{crrrc}
\toprule
$n$ & Type & Total & Missing & Rate \\
\midrule
7 & Odd & 22 & 1 & 4.5\%\\
8 & Even & 57 & 0 & 0.0\%\\
9 & Odd & 51 & 1 & 2.0\%\\
10 & Even & 156 & 0 & 0.0\%\\
11 & Odd & 158 & 6 & 3.8\%\\
12 & Even & 566 & 8 & 1.4\%\\
13 & Odd & 499 & 46 & 9.2\%\\
14 & Even & 1366 & 11 & 0.8\%\\
15 & Odd & 3978 & 354 & 8.9\%\\
16 & Even & 5900 & 103 & 1.7\%\\
17 & Odd & 7094 & 357 & 5.0\%\\
18 & Even & 19204 & 345 & 1.8\%\\
19 & Odd & 32577 & 2363 & 7.3\%\\
20 & Even & 118057 & 2297 & 1.9\%\\
21 & Odd & 2426 & 190 & 7.8\%\\
22 & Even & 1275 & 21 & 1.6\%\\
23 & Odd & 4003 & 234 & 5.8\%\\
24 & Even & 2920 & 54 & 1.8\%\\
25 & Odd & 9040 & 561 & 6.2\%\\
26 & Even & 4949 & 106 & 2.1\%\\
27 & Odd & 17385 & 777 & 4.5\%\\
28 & Even & 12203 & 306 & 2.5\%\\
29 & Odd & 44890 & 2136 & 4.8\%\\
30 & Even & 24925 & 534 & 2.1\%\\
31 & Odd & 72 & 1 & 1.4\%\\
32 & Even & 175 & 0 & 0.0\%\\
33 & Odd & 14 & 0 & 0.0\%\\
34 & Even & 172 & 0 & 0.0\%\\
35 & Odd & 24 & 0 & 0.0\%\\
36 & Even & 282 & 0 & 0.0\%\\
37 & Odd & 21 & 0 & 0.0\%\\
38 & Even & 338 & 0 & 0.0\%\\
39 & Odd & 33 & 0 & 0.0\%\\
40 & Even & 541 & 0 & 0.0\%\\
41 & Odd & 35 & 0 & 0.0\%\\
42 & Even & 747 & 0 & 0.0\%\\
43 & Odd & 63 & 0 & 0.0\%\\
44 & Even & 1017 & 0 & 0.0\%\\
45 & Odd & 106 & 0 & 0.0\%\\
47 & Odd & 105 & 0 & 0.0\%\\
49 & Odd & 196 & 0 & 0.0\%\\
51 & Odd & 264 & 0 & 0.0\%\\
53 & Odd & 377 & 0 & 0.0\%\\
\bottomrule
\end{tabular}
\end{table}

\begin{remark}
For $n\ge 21$ odd, all solutions possess nontrivial symmetry (primarily $180^\circ$ rotational symmetry, class \texttt{rot2}). No identity-class solutions are recorded in the database for $n\ge 21$; this likely reflects search limitations rather than a proven non-existence.
\end{remark}

\subsection{C$_4$ Rotational Symmetry Solutions}

Table~\ref{tab:c4counts} shows the exponential growth of C$_4$-symmetric (rot4) solutions, confirming that $D(n)=2n$ for all even $n$ is structurally guaranteed.

\begin{table}[H]
\centering
\caption{C$_4$ rot4 solution counts for even $n$ (Flammenkamp database, supplemented by mvr/no-three-in-line for $n\ge44$).}
\label{tab:c4counts}
\footnotesize
\begin{tabular}{crcr}
\toprule
$n$ & rot4 Solutions & $n$ & rot4 Solutions \\
\midrule
12 & 4 & 34 & 172\\
14 & 13 & 36 & 281\\
16 & 13 & 38 & 337\\
18 & 7 & 40 & 541\\
20 & 16 & 42 & 746\\
22 & 8 & 44 & 1,016\\
24 & 23 & 46 & 1,366\\
26 & 36 & 48 & 2,124\\
28 & 58 & 50 & 3,381\\
30 & 92 & 52 & 5,062\\
32 & 101 & 54--56 & (see text)\\
\bottomrule
\end{tabular}
\end{table}

For $n\ge44$, the count grows approximately $1.5\times$ per 2-step increment. The Flammenkamp database lists 7,696 rot4 solutions at $n=54$ and 10,441 at $n=56$ (mvr/no-three-in-line). No rot4 solution is known for $n=74$ as of the database cut date (2026-06-25).

% ============================================================
\section{The Even-$n$ Threshold}
\label{sec:even}

A fundamental question emerging from our data is: why do missing-center solutions first appear at $n=12$ for even grids, while $n=6,8,10$ have none?

\subsection{Distance Ring Capacity}

The number of distinct distance rings (squared Euclidean distances from center) grows with $n$:
\begin{itemize}
\item $n=6$: 6 rings (capacity 12 at $\le2$ pts/ring), needs 12 points.
\item $n=8$: 9 rings (capacity 18), needs 16.
\item $n=10$: 14 rings (capacity 28), needs 20.
\item $n=12$: 19 rings (capacity 38), needs 24.
\end{itemize}

The \emph{slack ratio} (capacity/need) increases with $n$, so ring capacity alone is not the bottleneck.

\subsection{Collinearity as the True Barrier}

We formalized the distance-ring constraints as a counting matrix $M[i][j]$ and proved that for $n=8$, the ring constraints alone admit continuous families of solutions. The fact that exhaustive search finds none implies that \textbf{collinearity constraints}, not ring capacity, prevent missing-center solutions at $n<12$.

The number of ring-pair interactions grows quadratically:
\begin{equation}
\binom{R}{2} = \frac{R(R-1)}{2},\quad R = \text{number of rings}
\end{equation}

At $n=12$, the 19 rings provide $171$ ring-pair interactions --- sufficient geometric diversity for both constraints to be satisfied simultaneously. The threshold is a genuine \emph{combinatorial phase transition}.

\begin{conjecture}
Within the catalogued $D_4$-inequivalent symmetry classes, missing-center solutions exist for all even $n$ with $12\le n\le 30$.
\end{conjecture}

\noindent\textbf{Limited window (even $n$).} The missing-center phenomenon for even grids is confined to a finite window: it first appears at $n=12$ and persists through $n=30$, but in the Flammenkamp database it vanishes again for every even $n\ge 32$ (the only symmetry classes recorded at those sizes are rot4 and dia2; rot4 solutions necessarily contain the center by the C$_4$ theorem, while dia2 solutions in the database also universally contain the center, though a formal proof for dia2 is not yet established). This mirrors the disappearance of odd-$n$ missing-center solutions at $n\ge 33$ (Section~\ref{sec:extinction}), and shows that missing-center solutions are a transient regime rather than a permanent feature of even grids.

\subsection{Diagonal Reflection and Ring-Sharing}

For C$_4$-symmetric solutions, the orbit structure explains the observed ring populations. The C$_4$ rotation $R$ generates orbits of size 4. When combined with the diagonal reflection $\sigma(i,j)=(j,i)$, two distinct C$_4$ orbits can coalesce on the same distance ring:

\begin{itemize}
\item Points on the diagonal ($i=j$ or $i=n-1-j$) give single orbits of 4 points.
\item Points off the diagonal give paired orbits of $4+4=8$ points.
\end{itemize}

This explains why all C$_4$ solutions have ring populations of only 4 or 8.

% ============================================================
\section{Odd-$n$ Analysis and the $n=11$ Phase Transition}
\label{sec:odd}

Odd-$n$ grids exhibit dramatically different behavior. The missing-center count grows from 1 (at $n=7,9$) to 2363 (at $n=19$), but the growth is not monotonic:

\begin{equation}
1 \xrightarrow{\times1} 1 \xrightarrow{\times6} 6 \xrightarrow{\times7.7} 46 \xrightarrow{\times7.7} 354 \xrightarrow{\times1.0} 357 \xrightarrow{\times6.6} 2363
\end{equation}

\subsection{The $n=15\to17$ Plateau}

The near-plateau at $n=15\to17$ ($354\to357$) is notable. A possible contributing factor is that $n=15=3\times5$ is both composite and $4k+3$, which may yield higher missing-center counts than its prime neighbours, pulling its value up to near the level of $n=17$. However, with only two data points at this transition, this remains a speculative observation. The three-factor model (Section~\ref{sec:regression}) attempts to quantify the effects of parity, primality, and ring-pair density, but the limited data ($n=7$--$20$) prevents a reliable test of compositeness as an independent factor.

\begin{conjecture}[Three-Factor Model]
\label{conj:three}
The missing-center count $M(n)$ for odd $n$ is governed by:
\[
M(n) \propto \exp(\alpha n) \cdot \beta^{\delta(n)} \cdot \gamma^{\rho(n)}
\]
where $\alpha$ is a base growth rate, $\beta>1$ is a compositeness boost factor, $\gamma>1$ is a $4k+3$ parity factor, $\delta(n)=1$ if $n$ is composite and $0$ otherwise, and $\rho(n)=1$ if $n\equiv3\pmod{4}$ and $0$ otherwise.
\end{conjecture}

\subsection{Why $n=11$ is the Threshold}

The transition at $n=11$ is driven by ring-pair collinearity density:
\begin{itemize}
\item $n=7$: 10 rings, 45 pairs --- sparse constraint graph.
\item $n=9$: 15 rings, 105 pairs --- manageable.
\item $n=11$: 20 rings, 190 pairs --- threshold crossed.
\item $n=19$: 51 rings, 1275 pairs --- extremely dense.
\end{itemize}

The number of ring pairs grows from 45 to 190 between $n=7$ and $n=11$, a $4.2\times$ increase that fundamentally changes the feasibility landscape.

% ============================================================
\section{C$_4$ Evolution Across Even $n$}
\label{sec:c4evol}

Analysis of all known rot4 solutions (for the five $n$ with detailed orbital data) reveals a remarkably clean structure:

\begin{table}[H]
\centering
\caption{C$_4$ orbital structure across even $n$. Ring counts shown as observed minimum--maximum where individual solutions differ; $n=72$ figures are from Heule's published record solution and were not independently re-decoded in this study.}
\label{tab:c4evol}
\begin{tabular}{crrcc}
\toprule
$n$ & Orbits & Rings & Orbits=$n/2$? & Ring Population \\
\midrule
12 & 6 & 6 & Yes & All 4 \\
14 & 7 & 6--7 & Yes & 4 or 8 \\
16 & 8 & 7--8 & Yes & All 4 \\
18 & 9 & 8--9 & Yes & 4 or 8 \\
72 & 36 & 34 & Yes & 4 or 8 \\
\bottomrule
\end{tabular}
\end{table}

\subsection{Structural Properties of C$_4$ Solutions}

The C$_4$ geometry imposes several structural properties on any solution.

\begin{theorem}[Orbit-Ring Invariance]
\label{thm:orbitring}
In any C$_4$-symmetric solution, all four points of a given C$_4$ orbit lie at the same squared Euclidean distance from the grid center.
\end{theorem}

\begin{proof}
The squared distance is $d(i,j) = (2i-(n-1))^2 + (2j-(n-1))^2$. A direct computation shows $d$ is invariant under the C$_4$ rotation $R(i,j) = (n-1-j,i)$:
\[
d(R(i,j)) = (2(n-1-j)-(n-1))^2 + (2i-(n-1))^2 = (-(2j-(n-1)))^2 + (2i-(n-1))^2 = d(i,j).
\]
Hence the four points $\{p, R(p), R^2(p), R^3(p)\}$ all share the same $d$ value.
\end{proof}

\begin{corollary}[Ring population is a multiple of~4]
\label{cor:ring4}
In any C$_4$-symmetric solution, the population of each distance ring is a multiple of $4$.
\end{corollary}

\begin{proof}
By Theorem~\ref{thm:orbitring}, each orbit contributes $4$ points to a single distance ring. The population of each ring is therefore $4$ times the number of selected orbits whose canonical cell has that $d$ value.
\end{proof}

\begin{corollary}[Orbit count]
\label{cor:orbitcount}
Every C$_4$-symmetric $2n$-point solution on an even $n\times n$ grid uses exactly $n/2$ distinct C$_4$ orbits.
\end{corollary}

\begin{proof}
Each orbit contributes $4$ points, and the solution has $2n$ points. Since orbits partition the point set, the number of orbits is $2n/4 = n/2$.
\end{proof}

\noindent\textbf{Observed ring population values.} In all known C$_4$ solutions (Flammenkamp database, $n=12\dots56$, and Heule's $n=72$), the multiplicity per ring---i.e., the number of selected orbits sharing a given $d^2$ value---is either $1$ or $2$, resulting in ring populations of $4$ or $8$. Rings with population $8$ arise when two orbits $(i,j)$ and $(j,i)$ (related by diagonal reflection) are both selected, because $d(i,j)=d(j,i)$. Whether larger multiplicities can occur in any valid C$_4$ solution is an open question.

% ============================================================
\section{Missing-Center Absence in Catalogued Symmetry Classes}
\label{sec:extinction}

\subsection{Even $n$}

C$_4$-symmetric solutions exist for every even $n\le 72$ (Flammenkamp database), with the count growing exponentially ($\sim1.5\times$ per 2-step). The structural argument via C$_4$ orbit theory strongly suggests such solutions exist at any even $n$, implying $D(n)=2n$ for all even $n$ is strongly supported by structural and computational evidence.

\subsection{Missing-Center Disappearance in Database Symmetry Classes for Odd $n\ge 33$}

The situation for odd $n$ is more nuanced. Heule's SAT solver~\cite{heule2026} established $D(65)=D(67)=D(69)=2n$ using rct4 (near-rot4) symmetry. However, $n=71$ remains unsolved --- the only gap $\le 72$.

\subsubsection{Absence of Missing-Center Solutions in Known Classes at $n\ge 33$}

We extended the analysis through $n=45$ using the Flammenkamp configuration database (catalogued classes continue to $n=72$). A dramatic structural phase transition occurs at $n=31$:

\begin{table}[H]
\centering
\caption{Missing-center disappearance in known database symmetry classes for odd $n\ge 31$.}
\label{tab:extinct}
\begin{tabular}{crrrcl}
\toprule
$n$ & Total & Missing & Rate & rct4 count & Symmetry \\
\midrule
31 & 72 & 1 & 1.4\% & 5 & dia1, dia2, rct4 \\
33 & 14 & \textbf{0} & \textbf{0\%} & 14 & rct4 \\
35 & 24 & 0 & 0\% & 23 & rct4, dia2 \\
37 & 21 & 0 & 0\% & 21 & rct4 \\
39 & 33 & 0 & 0\% & 33 & rct4 \\
41 & 35 & 0 & 0\% & 35 & rct4 \\
43 & 63 & 0 & 0\% & 63 & rct4 \\
45 & 106 & 0 & 0\% & 106 & rct4 \\
\bottomrule
\end{tabular}
\end{table}

As established in Theorem~\ref{thm:c2} (the C$_2$ Direction Theorem, Section~\ref{sec:c2}), the rot2 symmetry class requires all $n$ antipodal pairs to have distinct directions from center. This constraint, combined with the row-degree constraints, undergoes a sharp phase transition at $n=31$: $44{,}828$ solutions exist at $n=29$, yet zero at $n=31$ (with a single known survivor at $n=33$). The threshold is a structured combinatorial collapse whose precise algebraic cause remains open.

For $n\ge 35$, odd-$n$ solutions in the tracked database are predominantly rct4 (Table~\ref{tab:d4}), and \textbf{rct4 solutions invariably have the grid center as a circumcenter}. This is a group-theoretic necessity for the rct4 class: rct4 (and C$_4$) orbits on odd-$n$ grids have orbit size 4, forcing $\ge 4$ points per distance ring.

Thus, within the catalogued symmetry classes, missing-center solutions are absent for all odd $n\ge 33$. The center is always a circumcenter in every solution recorded in the Flammenkamp database above this threshold. \textbf{Caveat}: iden-class (non-symmetric) solutions are only tracked up to $n=20$ in the database, so the existence of iden-class missing-center solutions at larger odd $n$ remains an open question.

\noindent\textbf{Remark on the rct4 gap.} An interesting anomaly is that $n=11$, $n=13$, and $n=15$ have \emph{no} known rct4 solutions in the Flammenkamp database, while smaller ($n=9$) and larger ($n=17,19$) values do. This is not a theoretical impossibility but a search gap: the distance ring capacities for $m=5,6,7$ ($n=2m+1$) fall into a critical range where rct4's concentrated ring usage pattern cannot be satisfied. The gap could potentially be filled by a dedicated rct4-targeted search.

\subsubsection{D$_4$ Reconstruction Note}

The Flammenkamp database stores $D_4$-inequivalent solutions. Full solution counts are estimated by multiplying each class by its orbit size ($8\times$ iden, $4\times$ rot2/dia1, $2\times$ rot4, $1\times$ rct4), with cross-validation against exhaustive enumeration showing $\sim0.8\%$ agreement at $n=9$. A three-factor regression model (parity, primality, ring-pair density) fitted on $n=7$--$20$ achieves $R^2\approx0.83$ but with only 14 data points and 4 parameters, the risk of overfitting is substantial. Detailed regression results appear in Appendix~\ref{app:regression}.

Our analysis suggests that if $D(71)=2n$ exists, it would likely belong to the \textbf{rct4} symmetry class --- the only class known to support $2n$-point solutions for odd $n\ge 33$, as demonstrated by Heule's $n=65,67,69$ results~\cite{heule2026}. The rct4 class reduces the search to selecting fundamental-domain orbits and is the only viable structural class for large odd-$n$ grids.

\begin{conjecture}
$D(71)=2\times71=142$, and an rct4-symmetric solution exists. By the C$_4$ theorem, any such solution necessarily has the grid center as a circumcenter.
\end{conjecture}

% ============================================================
\section{Conclusion}
\label{sec:conc}

We have studied the geometric structure of extremal No-Three-In-Line configurations through the lens of symmetry-imposed constraints. Our main contributions are three theorems linking symmetry class to structural properties of maximal $2n$-point solutions:

\begin{enumerate}[nosep]
\item \textbf{C$_4$ Theorem} (Theorem~\ref{thm:c4}): $90^\circ$ rotational symmetry forces the grid center to be a circumcenter, because each C$_4$ orbit contributes $4$ points at the same distance.
\item \textbf{Row-Degree Theorem} (Theorem~\ref{thm:rowdeg}): For C$_4$-symmetric solutions on even $n=2m$, the row/column constraints are equivalent to a $2$-regular graph on $m$ vertices---a purely graph-theoretic characterization.
\item \textbf{C$_2$ Direction Theorem} (Theorem~\ref{thm:c2}): $180^\circ$-symmetric solutions require $n$ distinct center-to-point directions, with a sharp unsatisfiability transition at $n=31$ whose algebraic cause remains open.
\end{enumerate}

As a corollary, solutions that lack the center as a circumcenter---\emph{missing-center} configurations---are restricted to symmetry classes not covered by these theorems (primarily identity and $180^\circ$ rotation). Complete missing-center counts are provided for $n=2\dots13$ via exhaustive search and for $D_4$-inequivalent classes through $n=53$ via the Flammenkamp database. The data reveals a three-phase evolution: abundance ($n=7$--$19$), declining rate ($n=21$--$29$), and disappearance from catalogued classes ($n\ge 33$).

The remaining open case $n=71$ is the only gap $\le 72$ in the $D(n)=2n$ conjecture. Our structural analysis suggests that a solution, if it exists, would require rct4 symmetry and would necessarily have the grid center as a circumcenter.

% ============================================================
\section*{Open Problems}
\label{sec:open}

We conclude with three open problems that emerge naturally from this work.

\paragraph{1. The cycle structure conjecture.}
Theorem~\ref{thm:rowdeg} characterizes C$_4$-symmetric solutions as $2$-regular graphs on $m=n/2$ vertices. Among the $21{,}601$ known C$_4$ solutions for $n=12\dots56$, the two simplest cycle structures---a single $m$-cycle and a $(1,m-1)$ decomposition (one self-loop plus an $(m-1)$-cycle)---together account for $36.4\%$ of all solutions, and this fraction is stable across the entire range. Moreover, these two patterns are known to yield valid solutions for every tested $m\ge 10$.

\begin{conjecture}[Cycle Universality]
\label{conj:cycle}
For every even $n=2m\ge 20$, the $m$-cycle and $(1,m-1)$ cycle patterns produce C$_4$-symmetric $2n$-point No-Three-In-Line solutions. Equivalently, there exists a $2$-regular graph on $m$ vertices---specifically the cycle $C_m$ or the cycle-with-loop $C_{m-1}\cup L_1$---whose $4m$ realized points under the C$_4$ orbit map contain no three collinear.
\end{conjecture}

A proof of Conjecture~\ref{conj:cycle} would provide an explicit infinite family of $2n$-point solutions, establishing $D(2m)=2m$ constructively for all $m$.

\paragraph{2. The rot2 unsatisfiability proof.}
Theorem~\ref{thm:c2} reduces the C$_2$ existence problem to a direction-packing problem. The collapse at $n=31$---from $44{,}828$ solutions to zero---is sharp and unexplained. The constraint density crosses $74$ triples per available pair at the threshold, but the constraint graph is structured rather than random, so a random-$k$-SAT analogy is only qualitative. An algebraic or number-theoretic proof of the $n=31$ threshold (or the existence of a single obstruction invariant) would be a significant advance.

\paragraph{3. Missing-center solutions: finite or infinite?}
The missing-center property is destroyed by C$_4$ and rct4 symmetry, and the rot2 class becomes UNSAT at $n=31$. However, identity-class (non-symmetric) solutions---the largest source of missing-center configurations at small $n$---are only tracked to $n=20$ in public databases, and partial data reveals at least $17$ identity-class missing-center solutions at $n=21$. Whether missing-center solutions exist for arbitrarily large $n$, or whether every sufficiently large $2n$-point solution must have the grid center as a circumcenter, remains the most fundamental open question raised by this work.

All code and data are publicly available at \url{https://github.com/aujurd22/no3inline-missing-center} (version v1.0).

% ============================================================
\appendix
\section{Regression Models for Missing-Center Rates}
\label{app:regression}

\subsection{Reconstruction from D$_4$-Inequivalent Counts}

The Flammenkamp database stores $D_4$-inequivalent solutions. The full solution count is estimated by multiplying each class by its orbit size ($8\times$ iden, $4\times$ rot2/dia1, $2\times$ rot4, $1\times$ rct4). Cross-validation against exhaustive C++ enumeration ($n\le13$) shows agreement within $0.8\%$ at $n=9$.

\begin{table}[H]
\centering
\caption{Reconstructed full missing-center counts.}
\label{tab:fullrecon}
\footnotesize
\begin{tabular}{crrcc}
\toprule
$n$ & Full Total & Full Missing & Rate & Type \\
\midrule
7 & 132 & 4 & 3.03\% & Odd \\
8 & 380 & 0 & 0.00\% & Even \\
9 & 368 & 8 & 2.17\% & Odd \\
10 & 1,135 & 0 & 0.00\% & Even \\
11 & 1,120 & 36 & 3.21\% & Odd \\
12 & 4,348 & 52 & 1.20\% & Even \\
13 & 3,622 & 292 & 8.06\% & Odd \\
14 & $\sim$10,568 & $\sim$85 & 0.80\% & Even \\
15 & $\sim$30,634 & $\sim$2,726 & 8.90\% & Odd \\
16 & $\sim$46,304 & $\sim$808 & 1.74\% & Even \\
17 & $\sim$55,573 & $\sim$2,797 & 5.03\% & Odd \\
18 & $\sim$152,210 & $\sim$2,735 & 1.80\% & Even \\
19 & $\sim$258,170 & $\sim$18,729 & 7.26\% & Odd \\
20 & $\sim$941,580 & $\sim$18,326 & 1.95\% & Even \\
\bottomrule
\end{tabular}
\end{table}

\subsection{Three-Factor Model}

Using the $D_4$-inequivalent counts directly with weighted least squares:
\[
\text{missing rate} = -1.22 + 0.77\cdot(\text{mod4}) + 2.95\cdot(\text{prime}) + 0.44\cdot\log C(R,2),
\]
with $R^2 = 0.834$ (adjusted $\approx 0.78$, leave-one-out $Q^2 \approx 0.71$). Fitted on $n=7$--$20$ (14 points, 4 parameters); overfitting risk is substantial.

\subsection{Sum-of-Two-Squares Refinement}

A refined model incorporating the representation count $r_2(d)$:
\[
\text{missing rate} = 2.57 - 0.64\cdot(\text{max }4k+1\text{ primes}) + 0.12\cdot(\text{rings}) - 0.23\cdot(\text{r}_2=0)
\]
achieves $R^2=0.897$ ($Q^2\approx0.74$). The $4k+1$ prime coefficient flips sign under minor data changes, indicating fragility; ring count (a proxy for $n$) and primality are the only robust predictors.

% ============================================================
\bibliographystyle{plain}
\bibliography{references}

\end{document}
